\section*{2. Related work}
\paragraph{Collaborative learning for biological data.}
Parameter aggregation enables institutions to train a shared model while retaining their measurements locally \citep{McMahan2017,Sheller2020}. In single-cell analysis, scFed addresses cell-type classification and FedscGen addresses distributed batch correction \citep{Wang2024scFed,Bakhtiari2025FedscGen}. Tabula and its Chiron implementation support collaborative single-cell foundation-model training through an agent-enabled platform \citep{Ding2026Tabula,ChironCode}. Automated design can interact with local training: FedEx tunes hyperparameters while training shared weights \citep{Khodak2021FedEx}, and Helmsman synthesizes collaborative-learning systems through planning, coding and simulation-based refinement \citep{Li2025Helmsman}.

\paragraph{Self-improving agents in biological research.}
Self-improving agents turn experiment outcomes into subsequent computational proposals. AIDE searches candidate programs through drafting, debugging and improvement \citep{Jiang2025AIDE}, while the AI Scientist and AI-Scientist-v2 connect ideas to executed experiments, replanning, tree search and writing \citep{Lu2024AIScientist,Yamada2025AIScientistV2}. DrugEvolve combines research, implementation, evaluation and algorithm memory for drug-development tasks \citep{Zhou2026DrugEvolve,DrugEvolveCode}. BioCoLoop studies how multiple data-owning laboratories can supply development feedback through a shared research interface. TAPB, a ProteinTalks-derived efficacy head and scDEBART-based cell representations provide its task-specific prediction components \citep{lin2025tapb,Sun2026ProteinTalks,sung2026scdebart}. The proposal schema, fitting service and selection mechanism are common across tasks. Strong simple references on some perturbation benchmarks motivate evaluating the research procedure with fixed task predictors \citep{ahlmann2025linear}. Configuration search also connects to random search \citep{bergstra2012random}, Bayesian optimization \citep{snoek2012practical} and language-model proposal distributions \citep{yang2024opro}. We use feedback-free direct search to isolate the effect of development history under a shared proposal and fitting interface.

BioCoLoop connects collaborative fitting to a persistent development history. A shared coordinator executes candidate configurations across laboratories, aggregates their development measurements and returns structured evidence to a proposal model whose weights remain fixed. This interface supports history-guided proposal revision (Section 3.3) and, in a separate study, trajectory-guided training allocation (Section 3.4).
